% =============================================================================
% Chapter 7: Beyond the Far Field: Throat and Particle Program
% Input-ready fragment for the lean toy-model research proposal.
% This file intentionally contains no preamble or document environment.
% =============================================================================

\chapter{Beyond the Far Field: Throat and Particle Program}
\label{ch:throat-particle-program}

The common-medium phase in \cref{ch:cross-sector-integration} ends, at best,
with a stable finite brane, a complete background response, and a realized
far-field region for one frozen Candidate Parent System. That is not yet a
particle theory. The localized sources used in the far-field calculations still
have to be replaced by solutions of the parent equations whose geometry,
internal support, transport, conversion, couplings, response, and reservoir
loading all arise from the same candidate medium.

This chapter outlines that nonlinear program. It begins with one supported
throat, continues through species and orientation partners, constructs moving
and accelerating families, and then studies two-throat interactions from the
far field through finite-thickness contact. The final stage closes the feedback
among each local one-way throat drain, the global bulk reservoir, and the
separate distributed \(S_{\mathrm{leak}}\) return across the porous brane. No
new particle-only mechanism is introduced at this stage. The purpose is to
determine whether the structures permitted by Candidate V1 actually realize the
source classes, signs, normalizations, stability, and conservation assumptions
used in the earlier chapters.

\section{One-throat existence, support, and stability}

A candidate particle is a finite nonlinear throat through the ordered slab, not
a point source placed on an otherwise fixed background. Its solution must include
the full parent fields in \((\mathbf x,w)\), the deformed interfaces where those
interfaces remain meaningful, the surrounding density and flow, the local order
conversion region, the trapped support field, and the asymptotic reservoir data.
The apparent mouth seen by a brane observer is only one projection of this
complete object.

The first calculation is a spectral seed on a provisional throat geometry. The
transverse support field must solve the same variable-coefficient shear operator
that governs the light sector, now on a nonuniform defect background. Schematically,

\[
\mathcal L_T[\chi_B,n,\text{geometry}]\,\phi_T
=
\omega^2\mathcal W_T[\chi_B,n] \,\phi_T.
\]

A true bound seed requires finite weighted norm, finite energy, and acceptable
asymptotic decay. A leaky seed is a different problem: it requires outgoing
conditions, a complex pole or equivalent leakage rate, and a derived lifetime.
The support region may emerge as a rim, collar, sheath, cavity wall, or another
localized pattern of stress. It must be diagnosed from the solved energy and
stress distribution rather than imposed as a hard domain in advance. A
background light mode that is healthy on the homogeneous slab is not thereby a
bound throat mode; the defect spectrum, bulk continua, longitudinal coupling,
and order response must be recalculated.

The spectral seed is not the particle. It must be continued to a
self-consistent free-boundary or time-periodic solution. In that solution the
support mode's cycle-averaged stress holds open the aperture whose shape sets the
mode spectrum. The aperture supplies a localized one-way drain, and its
ordered-to-bulk conversion loads the global reservoir. Bulk material does not
reverse course through the throat. It returns separately by reorganizing into
brane state across the porous brane, through the positive bulk-to-brane
contribution to \(S_{\mathrm{leak}}\). Mechanical backpressure can still couple
the reservoir to the local throat boundary data without constituting material
return through the throat. The intended closure is therefore

\[
\boxed{
\begin{aligned}
E_{\rm support},\omega_{\rm standing}
&\longleftrightarrow R_{\rm throat},\text{shape}
\longrightarrow
\{Q_n^{\rm net},Q_n,Q_\chi\}_{\rm localized\ drain/conversion}
\longrightarrow
\{p_{\rm bulk},\mathcal B_\infty\},
\\[0.3em]
\{p_{\rm bulk},\mathcal B_\infty\}
&\longrightarrow
S_{\mathrm{leak}}^{\mathrm{bulk}\to\mathrm{brane}}(\mathbf x,t)
\longrightarrow
\text{distributed reordering across the porous brane}
\\
&\longrightarrow
\text{brane background and local throat boundary data}.
\end{aligned}
}
\]

Here the signed net constituent flux, the selected positive throughput on an
outward-drain branch, and the gross local ordered-to-bulk conversion diagnostic
remain distinct. Their relationship is an output of the chosen conservative,
relaxational, or explicitly mixed order dynamics. A critical-flow estimate may
supply a necessary throughput bound for a fixed mouth and upstream state, but it
cannot by itself select the throat radius, establish stability, or turn support
energy into observed mass.

The nonlinear continuation must determine the equilibrium or cycle-averaged
radius and shape, whether the core is anchored mainly to one interface or spans
the slab, the support frequency and localization, the transport and conversion
profiles, and every source component exported to the far field. In particular,
it must replace the provisional electric and gravity sources with the throat's
complete odd and even source vectors, its leading coupling to the Coulomb carrier,
and its local source-side transport, pressure, stress, and ordered-to-bulk
conversion data. The associated distributed porous-brane
\(S_{\mathrm{leak}}\) return profile is a separate output of the coupled global
problem. The gravity field must remain probe independent, and the same one-time
calibration used in \cref{ch:far-field-force-program} must be retained.

A converged stationary geometry may still contain a dynamically oscillating
support mode. Stability must therefore be tested about the complete solution,
not inferred from the seed eigenvalue. A conservative periodic throat requires a
Floquet or equivalent Hamiltonian response analysis. A relaxational or mixed
throat requires the full retarded linearization, including every internal
reservoir variable used to close damping and entropy production. The audit must
include translational zero modes, radius and shape oscillations, support-mode
perturbations, orientation-changing channels, negative-energy or growing modes,
and leakage into transverse, longitudinal, odd, even, bulk, and conversion
branches.

The one-throat stage passes only if at least one branch has a finite supported
geometry, nonzero branch-consistent throughput or conversion where required,
acceptable localization or lifetime, no unphysical instability, and source
profiles compatible with the already frozen far-field contracts. Failure of a
particular seed, support family, or order-dynamics branch rejects that route. If
no allowed branch of Candidate V1 can produce a sufficiently stable supported
throat, V1 has a common far-field medium but no particle realization.

\section{Species structure, orientation partners, and charge universality}

A solved throat family is labeled

\[
\mathcal T_{\alpha,s},
\qquad
s\in\{+1,-1\},
\]

where \(s\) is the normal orientation that supplies the proposed electric sign
and \(\alpha\) contains the orientation-independent species data. Those data may
include the support-mode family, core topology and boundary conditions,
equilibrium size and shape, transport and conversion structure, circulation or
chirality, and any additional conserved or metastable internal configuration.
The program must discover which such branches exist; it may not identify species
with orientation by definition.

Both orientations of each species must be solved from the same parent equations
under reflected environmental data. The complete antiparticle candidate is

\[
\boxed{
\mathcal C_{\rm th}
=
\mathcal I_{\rm internal}\circ\mathcal R_w
},
\]

not geometric reflection alone unless the solved branch shows that the internal
operation is trivial. Every orientation-odd phase, circulation, chirality,
support convention, topological datum, or conversion current must transform
consistently, while the required orientation-even scalars remain equal. For an
equilibrium branch this includes equal energy under reflected data. For a
conservative periodic branch it also includes a consistent map of support-mode
action, frequency, and phase and comparison in one declared cycle-averaged
ensemble. For a driven or mixed branch, equality of throughput, dissipation,
entropy production, lifetime, stationary stress and momentum fluxes, and
retarded response may replace or supplement an equilibrium energy test.

Consequently, reversing an electron-like branch produces the leading
positron-like candidate, not a proton-like object. A proton-like object would
have to occupy another internal branch, carry additional topology, or be a
composite throat structure. This distinction is essential both for annihilation
and for neutral binding: opposite electric signs do not imply identical
orientation-independent cores.

The same solved families must decide whether charge is universal and discrete.
If the far-field odd channel factorizes, its one-throat charge may be written

\[
q_{\alpha,s}=s\,g_\alpha,
\qquad
g_\alpha\ge0.
\]

The orientation supplies all sign, while \(g_\alpha\) is a derived residue or
coupling magnitude. Ordinary unit-charge species must produce the same leading
magnitude to the required accuracy even when their internal structures differ.
A common slab thickness is not enough: the core boundary data must be quantized,
topologically controlled, selected by a universal branch condition, or
sufficiently decoupled from the long-range carrier. Continuous moduli that allow
arbitrary charge, radius, throughput, or source strength would expose a major
particle-sector problem unless the resulting variation is part of an explicitly
bounded analog regime.

A binary label also does not establish additive conservation. The throat family
must support a conserved oriented count or current, with pair creation and
annihilation changing the local population while preserving the net permitted
quantity. Only after that construction and unit-coupling universality are derived
can a conventional long-range charge such as \(Q_{\rm eff}=g_0N_s\) be assigned.
The species study therefore requires more than one attractive solution: it must
continue several internal branches, measure their source residues and response
coefficients, identify any discrete invariants, and determine whether the
antiparticle map and charge law survive across the family without retuning.

\section{Moving throats, inertia, passive response, and radiation}

A stationary throat is next continued into a family with small uniform velocity,
external forcing, and controlled acceleration. The moving object is not a rigid
mouth translated through a passive medium. Its standing mode must retime and
rephase; its throat shape, near-field flow, pressure, density, order profile,
and local conversion region must be rebuilt, together with the surrounding
distributed porous-brane \(S_{\mathrm{leak}}\) return disturbance. The response
of that complete dressed configuration is the proposed origin of inertia and the
source of several observable effects that must remain separately identifiable.

For a closed reversible branch, the low-velocity inertial tensor can be extracted
from the total momentum of the complete solution,

\[
M_{ij}^{\rm inertial}
=
\left.\frac{\partial P_i}{\partial v_j}\right|_{\mathbf v=0},
\]

where \(P_i\) includes the support field, geometry, near medium, order sector,
and any reversible reservoir momentum. A relaxational or mixed branch instead
requires a frequency-dependent response of the form

\[
F_i(\omega)
=
\left[-\omega^2M_{ij}(\omega)-i\omega\Gamma_{ij}(\omega)+\cdots\right]
\delta x_j(\omega),
\]

with the damping, heat, entropy, and internal reservoir variables included in
the same closed system. The ordinary branch must have a positive symmetric
quasistatic inertial response. Its retarded response must also be passive for an
unexcited reservoir: a small periodic perturbation may store or dissipate work,
but it may not draw unlimited net power from an unspecified background.

Passive gravitational response is a different calculation. The source-side
gravity field and its calibration are fixed before the test throat is chosen.
The moving-solution family must then determine the low-frequency coefficient
that maps that external field into force or acceleration. Positive response is a
separate gate, followed by the question of whether the low-energy ratios among
local active source strength, passive response, and inertial response are
sufficiently universal across species. A calibration that absorbs
species-dependent response would not establish universality.

Motion also separates three effects that can otherwise be confused. The
sign-dependent transverse response proportional to the oriented current is the
electromagnetic magnetic sector. Sign-independent rebuilding of the surrounding
medium contributes to inertia or wake formation. A speed-dependent change in
throughput, conversion, or the signed gravity source is a preferred-medium
gravity effect. All three may occur, but they must be extracted from different
parity and response components of the same solution.

Acceleration then tests dynamic electromagnetic closure. The same conserved
oriented current and the same leading coupling that produced the static electric
residue must normalize the \(O(v)\) magnetic response and the \(O(a)\) transfer
of energy and momentum into both transverse light polarizations. The full
far-zone calculation must simultaneously measure

\[
P_{u_T},\quad
P_{\rm odd},\quad
P_{\rm even},\quad
P_{u_L},\quad
P_{\rm bulk},\quad
P_\chi,
\]

rather than reporting only the desired transverse radiation. Radiation into the
odd electric branch, even thickness response, longitudinal mode, bulk sound, or
order and reservoir channels is a prediction of Candidate V1 and may constrain
particle lifetime, acceleration response, and electromagnetic accuracy.

Uniform motion supplies an equally important no-drag test. In a reversible
branch, every mode coupled to the throat contributes to the directional
Landau-like emission threshold; a slow quadratic branch may drive the critical
velocity to zero unless its coupling vanishes by symmetry. In a dissipative or
mixed branch, the absence of an on-shell wave is not enough: low-frequency
friction, order-relaxation lag, diffusion, and conversion drag must be calculated
directly. The moving-throat stage succeeds only if one family provides positive
ordinary response, acceptable low-energy universality, one electric--magnetic--
radiative normalization, and bounded extra radiation and drag in the declared
regime.

\section{Two-throat interactions and finite-thickness crossover}

The two-body program begins from two independently solved throats, not from two
adjustable point sources. At large separation their measured source profiles
must reproduce the realized far-field forces of Candidate V1. The solutions are
then continued inward while tracking how each throat changes the other's core,
support mode, transport, conversion, distributed porous-brane return halo, and
reservoir data.
This continuation tests whether far-field factorization and approximate
additivity remain valid when the sources are no longer isolated.

Three separation regimes organize the calculation:

\begin{center}
\begin{tabularx}{\textwidth}{@{}L{0.22\textwidth}L{0.31\textwidth}X@{}}
\toprule
\textbf{Regime} & \textbf{Dominant description} & \textbf{Required output} \\
\midrule
\(R\gg H_{\rm br},R_{\rm throat}\)
& Normalized far-field modes and solved one-throat residues
& Recovery of gravity, odd electric, magnetic, and any allowed subleading forces with the frozen normalizations. \\
\(R\sim H_{\rm br}\)
& Full finite-thickness interface and parent-field response
& Cross-interface form factors, even thickness forces, polarization, overlap of porous-brane return profiles, and departure from point-source factorization. \\
\(R\) comparable to the nonlinear core scale
& Coupled parent-field two-core solution
& Core deformation, support-mode phase and stability, topology, contact, short-range force, radiation, and possible merger or exclusion. \\
\bottomrule
\end{tabularx}
\end{center}

The pair classes must include same-species same-orientation throats,
same-species opposite-orientation partners, and oppositely charged throats of
different or composite species. The odd Coulomb channel, reflection-even
thickness and material channels, gravity and distributed porous-brane return
response, magnetic terms, and nonlinear core interaction must remain separately
visible. A correction to
the odd monopole coefficient cannot be used to hide an even or core force, and
distributed porous-brane return-halo overlap must be allowed to change the signed
effective gravity source rather than being absorbed into an absolute mass
parameter.

The force construction depends on the solved branch. A genuine equilibrium pair
may admit an ensemble-appropriate static energy. A conservative periodic pair
may require fixed support-mode action, frequency, and relative phase or a
cycle-averaged stress force. A maintained dissipative pair requires the complete
stress, momentum-flux, conversion, and reservoir response unless a conservative
potential is independently established. This proposal does not reproduce the
full ensemble and Green-matrix program; those calculations remain in the
canonical ledger and the focused opposite-orientation companion. The proposal-
level requirement is that the physical force and its sign be derived from one
well-defined pair problem rather than selected by changing the ensemble.

A successful crossover calculation connects the far-field observables to a
regular finite-thickness interaction and identifies where the two-interface
graph description ceases to be adequate. It must also reveal environmental and
many-body corrections to charge, gravity, mass response, and source additivity.
An unphysical singularity, arbitrary branch jump, loss of conservation, or need
for a new short-range coefficient would reject the selected pair realization.

\section{Projected coincidence, core contact, and annihilation}

A brane observer records a projected position \(\mathbf x_a\), whereas the
nonlinear core occupies a finite region of the full space
\(\mathbf X=(\mathbf x,w)\). Oppositely oriented throats may therefore have the
same projected position while remaining separated across the slab. This is not
a small correction to point-particle language; it changes what must be calculated
before annihilation can even be defined.

The central logical separation is

\[
\boxed{
\text{odd electric cancellation}
\not\Rightarrow
\text{even material cancellation}
\not\Rightarrow
\text{core contact}
\not\Rightarrow
\text{annihilation}
}.
\]

For a complete same-species orientation pair, the leading odd source may cancel
at projected coincidence while the even thickness disturbance, support fields,
flow, conversion zones, and other orientation-independent structure add. If the
throats are periodic, exact instantaneous cancellation may additionally depend
on the relative support-mode phase supplied by the complete antiparticle map;
cycle-averaged electric neutrality does not remove possible phase-dependent
stress or sideband radiation.

While both slab interfaces remain graphs, a useful diagnostic separation is

\[
d_w^{\rm core}
=
H_{\rm br}+2h_t^{\rm tot}
-\ell_+^{\rm in}-\ell_-^{\rm in},
\]

where the even thickness disturbance changes the local interface spacing and
the inward reaches \(\ell_\pm^{\rm in}\) must come from the solved cores. A
positive value permits projected overlap with separated cores; zero marks first
contact within that description; a negative value signals that the separate-core
graph ansatz has failed. The formula is a diagnostic, not a substitute for the
nonlinear solution.

Once overhangs, necks, tubes, reconnection, pinch-off, or topology change occur,
the interface variables must give way to the complete parent fields. Contact
then requires a physical localized core diagnostic and a continuation rule that
keeps the two cores identifiable from large separation until their closures
first touch or their component topology changes. The result must be robust under
reasonable variation of the diagnostic threshold. After merger, separate
``plus'' and ``minus'' cores may not be reconstructed by an analyst-chosen split
of one total field. The detailed set definitions and numerical contact protocol
remain in the opposite-orientation companion.

Contact still does not select the outcome. A same-species throat--antithroat pair
may scatter, form a metastable neutral state, reconnect, annihilate into allowed
propagating and reservoir modes, or leave a neutral remnant. Because its net
oriented count can vanish, annihilation may be topologically permitted, but the
dynamics and conservation laws must show an actual channel. By contrast, two
oppositely charged objects of different internal or composite species may retain
unmatched odd multipoles, even material structure, support modes, or topology.
Electric neutrality must not force them into the same annihilation channel.

The decisive calculation is a full time-dependent collision evolution using the
parent fields. It must close constituent number, momentum, and energy and, for a
relaxational or mixed branch, entropy and internal reservoir accounting. It must
measure outgoing transverse light, odd and even disturbances, longitudinal and
bulk waves, conversion and reservoir excitation, and any surviving core. Failure
of immediate annihilation is not automatically a model failure; it may establish
a bound state, exclusion mechanism, or metastable channel. The particle--
antiparticle sector fails if no conservation-respecting annihilation or
neutralization channel exists for a complete same-species pair, or if the theory
forces materially different opposite charges to annihilate solely because their
leading electric monopoles cancel.

\section{Localized throat drainage, distributed porous-brane return, and the
reservoir fixed point}

Each throat supplies a localized, one-way drain. Material crossing its mouth
de-structures into bulk and cannot return through that throat. Return occurs
separately when bulk material reorganizes into brane state across the porous
brane, represented by the positive bulk-to-brane contribution to the signed
projected exchange term \(S_{\mathrm{leak}}\). The local drain, gross order
conversion, distributed return, and reservoir response therefore belong to one
conservation problem but are not one mechanism or one rate.

Every local throat calculation depends on asymptotic data supplied by the wider
medium. Denote the provisional environmental state by

\[
\mathcal B_\infty
=
\{n_\infty,P_\infty,\mu_\infty,\chi_\infty,T_\infty,
\text{distributed }S_{\mathrm{leak}}\text{ return profile},
\text{reservoir state}\},
\]

with only the variables appropriate to the selected dynamical branch retained.
A local solution obtained at fixed \(\mathcal B_\infty\) exports throughput,
conversion, stress, momentum, energy, radiation, and possible entropy to the
global bulk. Those outputs alter the reservoir and distributed return profile;
the changed global state then alters the backpressure and other local boundary
data without sending material back through the throat. The final closure is
therefore a fixed-point problem rather than a last one-way correction.

The iteration proceeds conceptually as follows:

\begin{enumerate}[leftmargin=2.2em,itemsep=0.36em]
  \item solve the supported local throat under provisional \(\mathcal B_\infty\);
  \item extract its constituent flux, order conversion, momentum and energy
        transfer, emitted disturbances, and branch-specific entropy production;
  \item solve the global bulk transport and reservoir evolution sourced by the
        population of throats, together with the spatially distributed
        bulk-to-brane return contribution to \(S_{\mathrm{leak}}\);
  \item update the backpressure, return scale, environmental asymmetry, gravity
        source profile, and local boundary data;
  \item repeat until the local and global solutions reproduce the same reservoir
        state, distributed \(S_{\mathrm{leak}}\) return profile, throat geometry,
        throughput, conversion, and emitted fluxes.
\end{enumerate}

The global problem must preserve the one-medium ontology. Drainage is not loss
into an external void. Material that crosses a throat remains in the common bulk
medium but does not return through that throat. Return is neither reverse throat
flow nor material supplied by an inexhaustible battery; it is the separate
distributed reorganization of bulk material into brane state across the porous
brane. A conservative branch must close total constituent number, momentum, and
energy. A relaxational or mixed branch must additionally track the internal heat,
disorder, entropy, and reservoir variables into which resolved order and
mechanical energy are transferred. Any maintained stationary throat or long-lived
particle must identify where the required energy resides and how the reservoir
changes over time.

The fixed point also completes several observables that cannot be closed locally.
The distributed porous-brane \(S_{\mathrm{leak}}\) return determines whether the
gravity source has a usable local matching region, becomes screened at a return
scale, overshoots, or changes sign.
The reservoir state influences throat radius, support frequency, lifetime,
conversion, drag, and the comparison of the two orientations. A nominally
reflection-symmetric particle branch may acquire an environmental splitting if
the two bulk-facing reservoirs are not symmetric; that effect must be calculated
rather than folded into the intrinsic antiparticle map.

A viable fixed point is stable or sufficiently long lived under perturbations of
both the local throat and the global reservoir. It must not require secular
reservoir depletion, unlimited entropy export to an unmodeled bath, or continual
retuning of boundary data. Several outcomes remain scientifically useful. A
stable fixed point would close the particle and reservoir program for the tested
regime. A slowly evolving solution could define a metastable particle sector
with a calculated lifetime and environmental dependence. A failure to converge
may reject the selected order-dynamics or distributed-return branch even if its
isolated throat looks acceptable. If every allowed branch of the frozen candidate lacks
a conservation-respecting local--global fixed point, Candidate V1 fails at the
final closure stage.

The result of this chapter is therefore not a promise that the far-field analog
automatically contains particles. It is a sequence of nonlinear existence,
response, interaction, topology, and reservoir tests that replaces each
provisional source with a solved object. Only a candidate that survives this
sequence can claim one connected classical analog extending from the common
brane and far field to stable particle-like throats and their global material
cycle.
